\chapter{Resultados y discusión}
\label{chap:resultados}


\section{Diseño de la evaluación}
\label{sec:diseno_evaluacion}

La evaluación se centra en comprobar la capacidad del generador para producir
modelos UVL con diferentes niveles de expresividad, la influencia de los
parámetros de configuración, su integración con Flamapy y UVLHub y el
funcionamiento de la batería de pruebas automatizadas. Para ello, se ejecutaron
pruebas funcionales, comprobaciones de compatibilidad con el parser UVL,
evaluaciones de integración y mediciones temporales bajo las configuraciones
descritas en las siguientes subsecciones.

La evaluación se realizó en el siguiente entorno:

\begin{itemize}
    \item \textbf{Sistema operativo:} Ubuntu 24.04.4 LTS.
    \item \textbf{Navegador:} Firefox 148.0.2 (64 bits).
    \item \textbf{Entorno de ejecución del generador:} Pyodide 0.27.7,
    utilizando Python 3.12.7 sobre WebAssembly.
    \item \textbf{Contenedores:} Docker 29.7.2, empleado para ejecutar
    UVLHub durante las pruebas de integración.
\end{itemize}

De forma adicional, el equipo anfitrión utilizado fue un Slimbook Executive con un procesador
Intel Core i7-13700H y 32 GB de memoria RAM.

\section{Resultados}
\label{sec:resultados_objetivos}


\subsection{Resultados sobre el objetivo O1}
\label{subsec:resultados_o1}


El objetivo O1 consistía en desarrollar un generador automático capaz de
producir modelos UVL con diferentes niveles de complejidad y expresividad,
incluyendo construcciones avanzadas del lenguaje que permitan ampliar la
variedad de modelos disponibles para el análisis de variabilidad.

Para verificar este objetivo se realizó una evaluación experimental basada en
tres configuraciones acumulativas de expresividad. En cada configuración se
generaron diez modelos utilizando la semilla fija
\texttt{20260922}. Todos los modelos se cargaron posteriormente mediante
\texttt{UVLReader} para comprobar su compatibilidad sintáctica con UVL.

La primera configuración incluía el nivel booleano y la cardinalidad de grupo,
pero no generaba atributos. Estos modelos se sometieron además a una
comprobación de satisfacibilidad mediante PySAT con posterioridad a la
generación, manteniendo desactivada la opción de generación condicionada por
satisfacibilidad. De este modo, la proporción observada de modelos
satisfacibles no queda alterada por reintentos del generador. Las
configuraciones aritméticas y tipadas se utilizaron para validar la generación y
el parseo de construcciones de mayor expresividad, pero no se incluyeron en la
evaluación SAT.

La Figura~\ref{fig:configuracion-O1} muestra la configuración de mayor
expresividad utilizada en las ejecuciones de validación sintáctica, con la
semilla \texttt{20260922}. No representa la ejecución booleana sometida a
SAT: esta se realizó con la generación de atributos desactivada y con
\textit{Ensure satisfiable} desactivado. La evidencia cuantitativa se
recoge en las Tablas~\ref{tab:validacion-o1} y
\ref{tab:validacion-o1-sat}.

\begin{figure}[H]
    \centering
    \includegraphics[width=0.40\textwidth]{configuracion_verificar_O1.png}
    \caption{Configuración usada para verificar O1. Captura realizada por el autor.}
    \label{fig:configuracion-O1}
\end{figure}

El Apéndice~\ref{app:modelos-o1} recoge tres ejemplos ilustrativos de modelos
generados y los parámetros utilizados en esas ejecuciones. Son ejecuciones
independientes de la evaluación cuantitativa de O1, que empleó la semilla
\texttt{20260922}. En particular, el modelo booleano del apéndice contiene un
atributo, mientras que la configuración utilizada para la comprobación SAT se
ejecutó sin atributos. Por tanto, los listados del apéndice no forman parte
del recuento de treinta modelos ni de los resultados SAT de las tablas.

La siguiente Tabla~\ref{tab:validacion-o1} resume la validación sintáctica mediante
\texttt{UVLReader}, mientras que la Tabla~\ref{tab:validacion-o1-sat} presenta
los resultados de la comprobación de satisfacibilidad. La columna ``UVLReader OK'' 
indica el número de
modelos cargados correctamente mediante el parser de Flamapy. Las columnas
relativas a SAT solo se aplican a la primera configuración, formada únicamente
por construcciones booleanas y sin atributos. ``SAT evaluados'' recoge los
modelos cuya comprobación terminó, ``SAT satisfacibles'' aquellos para los que
PySAT devolvió un resultado positivo, ``SAT no satisfacibles'' aquellos para
los que devolvió un resultado negativo y ``SAT no concluyentes'' los casos que
no finalizaron dentro del tiempo disponible.

\begin{table}[H]
\centering
\small
\begin{tabularx}{\textwidth}{@{}Xrrr@{}}
\toprule
\textbf{Configuración} &
\textbf{Modelos} &
\shortstack{\textbf{UVLReader}\\\textbf{OK}} &
\shortstack{\textbf{No}\\\textbf{parsean}} \\
\midrule
Booleano y cardinalidad de grupo, sin atributos
& 10 & 10 & 0 \\

Booleano, aritmética, cardinalidad de características y funciones agregadas
& 10 & 10 & 0 \\

Booleano, aritmética y tipos, con todos los niveles secundarios
& 10 & 10 & 0 \\

\midrule
\textbf{Total} & \textbf{30} & \textbf{30} & \textbf{0} \\
\bottomrule
\end{tabularx}
\caption{Validación sintáctica de los modelos generados para O1.}
\label{tab:validacion-o1}
\end{table}

\begin{table}[H]
\centering
\small
\begin{tabularx}{\textwidth}{@{}Xrrrrr@{}}
\toprule
\textbf{Configuración} &
\shortstack{\textbf{Modelos}\\\textbf{generados}} &
\shortstack{\textbf{SAT}\\\textbf{evaluados}} &
\shortstack{\textbf{Satisfacibles}} &
\shortstack{\textbf{No}\\\textbf{satisfacibles}} &
\shortstack{\textbf{No}\\\textbf{concluyentes}} \\
\midrule
Booleano y cardinalidad de grupo, sin atributos
& 10 & 9 & 1 & 8 & 1 \\
\bottomrule
\end{tabularx}
\caption{Comprobación SAT de la configuración booleana evaluada para O1.}
\label{tab:validacion-o1-sat}
\end{table}

El modelo restante de la configuración booleana se identificó como un caso no
concluyente, ya que la comprobación SAT no finalizó dentro del tiempo máximo
establecido para cada modelo (30 segundos). Por este motivo no se contabiliza como satisfacible 
ni como no satisfacible.

Los resultados muestran que los treinta modelos generados fueron aceptados por
\texttt{UVLReader}, por lo que las tres configuraciones produjeron artefactos
compatibles sintácticamente con el lenguaje UVL. En la configuración booleana,
la comprobación SAT terminó para nueve de los diez modelos: uno resultó
satisfacible y ocho no satisfacibles. Que un modelo generado sea no satisfacible
no implica que su generación sea incorrecta, ya que las restricciones aleatorias
pueden ser incompatibles entre sí; precisamente por ello la comprobación SAT
permite distinguir entre validez sintáctica y consistencia semántica.

La comprobación SAT se aplicó únicamente a la configuración booleana, por lo
que los resultados obtenidos no permiten establecer cómo varía la
satisfacibilidad al aumentar la complejidad de los modelos. Como consideración
teórica, un mayor número de restricciones o de características puede aumentar
el espacio de búsqueda del solucionador; sin embargo, su efecto sobre el
tiempo de resolución y la existencia de soluciones depende de la estructura
concreta del modelo. Para evaluar esta relación sería necesario realizar
experimentos específicos con distintas configuraciones y tamaños.

Por tanto, los resultados permiten verificar el cumplimiento del criterio de
éxito asociado al objetivo O1 en cuanto a la generación de modelos con
diferentes niveles de expresividad y a su compatibilidad sintáctica con UVL.
La comprobación SAT aporta además evidencia sobre la consistencia semántica de
la configuración booleana evaluada, aunque no se ha aplicado a las
configuraciones con construcciones aritméticas, tipadas y de cadenas.

\subsection{Resultados sobre el objetivo O2}
\label{subsec:resultados_o2}

El objetivo O2 consistía en desarrollar un sistema de configuración flexible
que permitiera modificar los parámetros principales del proceso de generación
y adaptar los modelos obtenidos a distintos escenarios.

Para verificar este objetivo se ejecutaron pruebas que comprueban tanto la
propagación de los parámetros a través del asistente de configuración como su
efecto sobre la estructura y el contenido de los modelos UVL generados. La
Tabla~\ref{tab:evidencias-o2} resume las comprobaciones realizadas.

\begin{table}[H]
\centering
\footnotesize
\setlength{\tabcolsep}{5pt}
\renewcommand{\arraystretch}{1.25}
\begin{tabularx}{\textwidth}{@{}
>{\raggedright\arraybackslash}p{3.1cm}
>{\raggedright\arraybackslash}p{7.0cm}
>{\raggedright\arraybackslash}X
@{}}
\toprule
\textbf{Funcionalidad} & \textbf{Prueba automatizada} &
\textbf{Comprobación realizada} \\
\midrule

Propagación de la configuración
& \path{test_step2_configuration_is_propagated_to_step3}
& Comprueba que las opciones seleccionadas en el segundo paso se trasladan correctamente al siguiente. \\

Cardinalidad de grupo
& \path{test_group_cardinality_off_produces_no_group_relations}

\path{test_group_cardinality_on_with_weight_produces_group_groups}
& Verifica la ausencia o generación de relaciones de cardinalidad de grupo. \\

Cardinalidad de características
& \path{test_feature_cardinality_off_produces_none}

\path{test_feature_cardinality_on_produces_some}
& Comprueba que la opción modifica la salida UVL generando o eliminando cardinalidades. \\

Niveles aritmético y de tipos
& \path{test_arithmetic_level_produces_arith_constraints}

\path{test_string_level_produces_len_constraints}
& Verifica la generación de expresiones aritméticas y restricciones mediante \texttt{len}. \\

Distribución de atributos
& \path{test_dist_boolean_only_produces_only_boolean_attrs}

\path{test_dist_integer_only_produces_integer_attrs}

\path{test_real_attribute_distribution_generates_real_attributes}
& Comprueba que las distribuciones configuradas determinan los tipos de atributos generados. \\

Cantidad de restricciones y variables
& \path{test_min_max_constraints_respected}

\path{test_min_vars_per_constraint_respected}
& Comprueba los límites aplicados durante la generación; el número
efectivo de restricciones puede ser inferior al mínimo configurado si
no hay variables compatibles para construirlas. \\

Semillas
& \path{test_determinism_same_seed_same_output}

\path{test_different_seed_different_output}
& Verifica la reproducibilidad con la misma semilla y la variación con semillas diferentes. \\

\bottomrule
\end{tabularx}
\caption{Pruebas técnicas del efecto de los parámetros de configuración.}
\label{tab:evidencias-o2}
\end{table}


Las pruebas anteriores no se limitan a comprobar que el asistente acepta los
valores introducidos, sino que verifican su efecto sobre los modelos
generados. Por ejemplo, al activar la cardinalidad de grupo se generan las
relaciones correspondientes, mientras que al activar el nivel aritmético
aparecen expresiones numéricas en las restricciones. De esta forma, existe una
relación verificable entre la configuración seleccionada, el comportamiento del
generador y el contenido del modelo UVL obtenido.

El criterio de éxito asociado al objetivo O2 se considera cumplido, ya que las
pruebas confirman tanto la correcta propagación de los parámetros por el flujo
de configuración como su influencia efectiva en los modelos generados.

\subsection{Resultados sobre el objetivo O3}
\label{subsec:resultados_o3}

El objetivo O3 consistía en integrar el generador dentro del ecosistema
Flamapy y permitir su utilización desde UVLHub como herramienta de generación
de modelos de características.

La integración se verifica mediante el flujo real de ejecución utilizado por
la aplicación. Los parámetros almacenados por UVLHub se convierten en una
instancia de \texttt{FmgeneratorModel} mediante
\texttt{from\_flat\_dict}. Posteriormente, el \textit{wrapper} importa directamente la operación
\texttt{GenerateFeatureModel}, crea una instancia y la ejecuta con el modelo
de configuración recibido. Como resultado se obtiene un
\texttt{FeatureModel} compatible con el metamodelo de Flamapy.

De forma independiente, la prueba \newline
\texttt{test\_generate\_feature\_model\_is\_discoverable\_by\_flamapy}
verifica que Flamapy también localiza esta operación mediante
\texttt{DiscoverMetamodels}. Así, se comprueba tanto su uso efectivo desde
UVLHub como su registro como operación descubrible dentro del ecosistema
Flamapy.

Esta integración se comprobó mediante las pruebas recogidas en la
Tabla~\ref{tab:evidencias-o3}.

\begin{table}[H]
\centering
\small
\setlength{\tabcolsep}{5pt}
\renewcommand{\arraystretch}{1.25}
\begin{tabularx}{\textwidth}{@{}
>{\raggedright\arraybackslash}p{6.4cm}
>{\raggedright\arraybackslash}X
@{}}
\toprule
\textbf{Prueba} & \textbf{Evidencia obtenida} \\
\midrule

{\footnotesize\path{test_wrapper_generates_feature_model_with_received_arguments}}
&
Comprueba que el \textit{wrapper} recibe los parámetros y genera un modelo de
características. \\

{\footnotesize\path{test_wrapper_build_one_model_delegates_to_feature_model_generator}}
&
Verifica que la generación se delega al componente
\texttt{GenerateFeatureModel}. \\

{\footnotesize\path{test_wrapper_serializes_uvl_and_reads_includes}}
&
Comprueba la serialización del modelo mediante Flamapy y la conservación de las
declaraciones \texttt{include}. \\

{\footnotesize\path{test_generate_feature_model_is_discoverable_by_flamapy}}
&
Comprueba que \texttt{DiscoverMetamodels} localiza la operación
\texttt{GenerateFeatureModel}, que esta puede ejecutarse y que devuelve un
\texttt{FeatureModel} válido. \\

\bottomrule
\end{tabularx}
\caption{Pruebas técnicas de integración con Flamapy y UVLHub.}
\label{tab:evidencias-o3}
\end{table}

La evidencia anterior demuestra que el componente no se limita a estar
presente en la estructura del proyecto, sino que es utilizado durante la
ejecución, produce un \texttt{FeatureModel} de Flamapy y genera modelos UVL
que pueden ser consumidos posteriormente por otras herramientas del
ecosistema.

Por tanto, el criterio de éxito asociado al objetivo O3 se considera cumplido.

\subsection{Resultados sobre el objetivo O4}
\label{subsec:resultados_o4}

El objetivo O4 consistía en evaluar el comportamiento temporal del generador
en escenarios con distintos volúmenes y tamaños de modelos de características.

Para ello se definieron tres perfiles de tamaño. En cada perfil se fijaron los
valores mínimo y máximo configurados para las características, los atributos y
las restricciones. El resto de parámetros, incluidos los niveles de
expresividad UVL y las distribuciones empleadas, se mantuvo constante. Estos
perfiles fijan los valores solicitados al generador, pero no garantizan que el
número de restricciones incorporadas sea siempre idéntico, ya que alguna
puede omitirse si no se consigue construir con las variables disponibles.

La Tabla~\ref{tab:perfiles-o4} recoge los valores configurados para cada
perfil.

\begin{table}[H]
\centering
\small
\renewcommand{\arraystretch}{1.2}
\begin{tabular}{lrrr}
\toprule
\textbf{Perfil} & \textbf{Características configuradas} &
\textbf{Atributos configurados} & \textbf{Restricciones configuradas} \\
\midrule
Pequeño & 10  & 5  & 10  \\
Mediano & 50  & 25 & 50  \\
Grande   & 100 & 50 & 100 \\
\bottomrule
\end{tabular}
\caption{Perfiles de tamaño utilizados en la evaluación temporal.}
\label{tab:perfiles-o4}
\end{table}

Para cada perfil se generaron lotes de 100, 500 y 1000 modelos UVL. Cada una
de las nueve combinaciones se ejecutó tres veces con la misma configuración y
semilla. El tiempo de generación se midió una vez que Pyodide y las
dependencias necesarias estaban disponibles, por lo que no incluye la carga
inicial del entorno ni el renderizado o la descarga de los resultados.

La carga de Pyodide y de las dependencias se midió de forma independiente en
diez ejecuciones. Esta separación permite distinguir el coste inicial de
preparación del entorno del coste propio de generar los modelos.

La Tabla~\ref{tab:carga-pyodide-o4} muestra la media y la desviación
típica del tiempo de preparación medidas en diez ejecuciones.

\begin{table}[H]
\centering
\small
\renewcommand{\arraystretch}{1.2}
\begin{tabular}{lrrr}
\toprule
\textbf{Fase} & \textbf{Repeticiones} & \textbf{Media (s)} &
\textbf{Desv. típica (s)} \\
\midrule
Carga de Pyodide y dependencias & 10 & 4,28 & 0,24 \\
\bottomrule
\end{tabular}
\caption{Tiempo de preparación del entorno de ejecución.}
\label{tab:carga-pyodide-o4}
\end{table}

La Tabla~\ref{tab:generacion-o4} recoge la media y la desviación típica
muestral de los tiempos de generación. Cada resultado se ha calculado a partir
de tres repeticiones.

\begin{table}[H]
\centering
\small
\renewcommand{\arraystretch}{1.2}
\begin{tabular}{lrrr}
\toprule
\textbf{Tamaño} & \textbf{Modelos} & \textbf{Media (s)} &
\textbf{Desv. típica (s)} \\
\midrule
Pequeño & 100  & 0,60 & 0,02 \\
Pequeño & 500  & 2,90 & 0,21 \\
Pequeño & 1000 & 5,86 & 0,23 \\
\midrule
Mediano & 100  & 1,89 & 0,16 \\
Mediano & 500  & 8,76 & 0,31 \\
Mediano & 1000 & 17,52 & 0,31 \\
\midrule
Grande & 100  & 3,77 & 0,04 \\
Grande & 500  & 18,47 & 0,30 \\
Grande & 1000 & 36,95 & 0,28 \\
\bottomrule
\end{tabular}
\caption{Tiempo de generación de modelos UVL según el tamaño del modelo y el
volumen del lote. Cada resultado corresponde a tres repeticiones.}
\label{tab:generacion-o4}
\end{table}

Los resultados muestran que, para un mismo perfil, el tiempo de generación
aumenta de forma aproximadamente proporcional al número de modelos solicitado.
Por ejemplo, en el perfil pequeño se pasó de 0,60 segundos para 100 modelos a
5,86 segundos para 1000 modelos. Este comportamiento se mantiene en los tres
tamaños evaluados.

El perfil estructural también se asocia a diferencias en el tiempo de
generación. Para lotes de 1000 modelos, la media aumentó de 5,86 segundos en el
perfil pequeño a 17,52 segundos en el mediano y 36,95 segundos en el grande.
Esta comparación refleja el efecto conjunto de los perfiles, ya que aumentan
simultáneamente las características, los atributos y las restricciones; no
permite atribuir el incremento a uno solo de estos parámetros.

La variabilidad entre repeticiones fue reducida. La mayor desviación típica
registrada fue de 0,31 segundos, obtenida en los lotes medianos de 500 y 1000
modelos. Los resultados son, por tanto, consistentes en el entorno evaluado.

La carga de Pyodide y de sus dependencias supuso un coste inicial medio de
4,28 segundos, independiente del tamaño y del número de modelos del lote. Una
vez preparado el entorno, los tiempos de la
Tabla~\ref{tab:generacion-o4} representan exclusivamente el coste de producir
los modelos UVL.

Por tanto, el objetivo O4 se considera cumplido, ya que se ha cuantificado el
comportamiento temporal del generador en nueve escenarios que combinan tres
volúmenes de generación y tres tamaños de modelo, diferenciando el coste de
preparación del entorno del proceso de generación.

\subsection{Resultados sobre el objetivo O5}
\label{subsec:resultados_o5}

El objetivo O5 consistía en validar el correcto funcionamiento del generador
mediante una batería de pruebas automatizadas que cubriese tanto la lógica
interna del motor de generación como el flujo de configuración integrado en
UVLHub.

El Capítulo~\ref{chap:pruebas} documenta los casos de las dos suites, sus
comandos de ejecución y su cobertura. La Tabla~\ref{tab:casos_por_fichero} recoge los casos superados por fichero, la Tabla~\ref{tab:cobertura_ficheros} muestra la cobertura de sentencias del código principal, y la Tabla~\ref{tab:resumen_suite} recoge
el resultado global de las pruebas. 
Las suites incluyen pruebas preexistentes y pruebas desarrolladas o adaptadas durante este
TFG, como se indica en la Sección~\ref{sec:estrategia_pruebas}.

El criterio de éxito asociado a O5 se considera cumplido para los escenarios
cubiertos por la batería automatizada: las dos suites finalizaron sin pruebas
fallidas. Este resultado permite valorar la validación realizada, sin extender
la conclusión a comportamientos que no formen parte de los casos probados.

\subsection{Resultados sobre los objetivos formativos O6--O10}
\label{subsec:resultados_objetivos_formativos}

Los objetivos O6--O10 se refieren a competencias adquiridas durante el
desarrollo del proyecto. A diferencia de los objetivos funcionales, no se
evalúan mediante una métrica de rendimiento, sino mediante evidencias
documentales y técnicas generadas durante el trabajo.

\begin{table}[H]
\centering
\small
\setlength{\tabcolsep}{5pt}
\renewcommand{\arraystretch}{1.2}
\begin{tabularx}{\textwidth}{@{}
>{\centering\arraybackslash}p{0.9cm}
>{\raggedright\arraybackslash}p{3.8cm}
>{\raggedright\arraybackslash}X
@{}}
\toprule
\textbf{Obj.} & \textbf{Evidencia} & \textbf{Evidencia documentada} \\
\midrule
O6 & Estudio del enfoque de generación propuesto por Sundermann et al. y
adaptación de sus bases conceptuales al componente implementado. &
El conocimiento científico revisado se ha aplicado en el diseño y desarrollo
del generador. \\

O7 & Integración de un componente nuevo en los proyectos Flamapy y UVLHub,
utilización de Docker y adaptación de pruebas existentes. &
El desarrollo se ha realizado sobre un ecosistema software real, con código
previo y componentes reutilizables. \\

O8 & Implementación en Python, adaptación de rutas y asistentes de Flask,
ejecución mediante Docker y desarrollo de pruebas automatizadas. &
Las tecnologías principales del proyecto se han utilizado de forma práctica
durante la construcción e integración de la solución. \\

O9 & Consulta de bibliografía científica, documentación técnica de Flamapy,
UVL y Pyodide, y justificación de decisiones de diseño. &
Las decisiones técnicas se han apoyado en información seleccionada y analizada
durante el desarrollo. \\

O10 & Planificación del trabajo, gestión de riesgos, documentación de diseño,
implementación, pruebas y resultados. &
El proyecto se ha gestionado y documentado de forma continuada durante sus
distintas fases. \\
\bottomrule
\end{tabularx}
\caption{Evidencias asociadas a los objetivos formativos. Elaboración propia.}
\label{tab:evidencias-objetivos-formativos}
\end{table}

La Tabla~\ref{tab:evidencias-objetivos-formativos} documenta la aplicación de
los objetivos formativos durante el desarrollo del proyecto. Estas evidencias
no permiten medir de forma cuantitativa el grado individual de adquisición de
cada competencia, por lo que los resultados deben interpretarse como evidencia
documental de su aplicación y no como una medición objetiva de su dominio.

\section{Comparación con trabajos previos}
\label{sec:comparacion_trabajos}



La Tabla~\ref{tab:comparativa_generadores} del Capítulo~\ref{chap:estado_arte}
recoge una comparación cualitativa entre diferentes herramientas relacionadas
con la generación automática de modelos de características. En ella se
consideran tres aspectos principales: la capacidad de generación automática, el
grado de personalización ofrecido y la incorporación de construcciones de
expresividad avanzada.

Los resultados obtenidos durante la evaluación del generador desarrollado
permiten observar que la solución propuesta cubre los tres criterios analizados.
A diferencia de S.P.L.O.T. y FeatureIDE, cuyas capacidades de generación están
principalmente orientadas a modelos basados en construcciones booleanas,
el generador desarrollado permite producir modelos UVL incorporando elementos
adicionales como atributos tipados, restricciones aritméticas, restricciones
sobre cadenas, funciones agregadas y diferentes mecanismos de cardinalidad.

En comparación con BeTTy, la solución desarrollada mantiene la capacidad de
configuración flexible de los parámetros de generación, pero amplía el rango de
expresividad de los modelos generados. Mientras que BeTTy permite definir
escenarios experimentales con diferentes niveles de complejidad estructural, el
generador desarrollado incorpora además construcciones propias de lenguajes de
variabilidad más recientes como UVL.

Además, la integración realizada dentro del ecosistema Flamapy y UVLHub supone
una diferencia respecto a las herramientas analizadas, ya que permite utilizar
el generador directamente dentro de una plataforma web orientada a la gestión y
análisis de modelos de características, manteniendo la compatibilidad con las
herramientas existentes del ecosistema.

No se ha realizado una comparación cuantitativa directa en términos de tiempo de
generación o volumen máximo de modelos, debido a que las herramientas estudiadas
no proporcionan escenarios de evaluación homogéneos que permitan establecer una
comparación justa, por lo que la comparación se ha centrado en las
capacidades funcionales y de expresividad.

En consecuencia, el trabajo desarrollado destaca por la generación de modelos
con mayor expresividad y por su integración dentro del ecosistema UVLHub y
Flamapy. No obstante, esta ventaja funcional no implica una superioridad global
frente a las herramientas analizadas, ya que algunas de ellas presentan una
mayor madurez, una validación experimental más extensa o métricas comparativas
que no se han reproducido en este trabajo.

\section{Discusión}
\label{sec:discusion}


Los resultados obtenidos muestran que la solución desarrollada responde al
problema planteado inicialmente: la necesidad de disponer de un generador capaz
de producir modelos de características con una mayor expresividad que las
aproximaciones tradicionales basadas principalmente en construcciones
booleanas. La integración del generador dentro del ecosistema Flamapy y UVLHub
permite disponer de una herramienta flexible para la generación de modelos UVL
destinados a tareas posteriores de análisis y experimentación.

Durante el desarrollo se encontraron diferentes dificultades relacionadas con
la integración y adaptación de los componentes existentes. Sin embargo, la evaluación confirmó el funcionamiento de las funcionalidades en los
escenarios y combinaciones cubiertos por las pruebas, incluyendo los niveles
de expresividad y los parámetros examinados.

La incorporación de construcciones avanzadas implica ciertos compromisos. Por
un lado, aumentar la expresividad mediante atributos tipados, restricciones
aritméticas, funciones agregadas o restricciones sobre cadenas incrementa la
complejidad interna del generador. Por otro lado, ofrecer una configuración más
flexible permite adaptar la generación a distintos escenarios, aunque aumenta
el número de parámetros que el usuario debe comprender. Asimismo, una mayor
riqueza en los modelos generados puede suponer un incremento del coste
computacional en configuraciones más complejas.

Finalmente, existen algunas amenazas a la validez de los resultados. La
evaluación confirma la compatibilidad sintáctica de los modelos generados con
UVL y con las transformaciones empleadas, pero no determina si dichos modelos
son representativos de dominios reales ni si resultan adecuados como conjunto de
evaluación para motores de razonamiento. Asimismo, los resultados temporales se
han obtenido con una única configuración de generación y en un entorno concreto,
por lo que podrían variar al modificar el tamaño o la complejidad de los
modelos. Tampoco se ha realizado una comparación cuantitativa directa con otras
herramientas bajo condiciones experimentales equivalentes.


Este capítulo ha contrastado la evidencia disponible con los criterios de
éxito y ha discutido el alcance de los resultados obtenidos. El
Capítulo~\ref{chap:conclusiones} sintetiza la aportación del trabajo, sus
limitaciones y las líneas de desarrollo que se derivan de ellas.
